\chapter{系统软件设计思想、框架与流程}
\section{分层设计与职责分配}
软件采用“硬件协议服务—应用控制—界面与存储”分层思路。所有板端请求由一个 \code{MonitorService} 管理，界面、音乐与游戏通过消息和队列协作。FPGA负责时间确定的采集与缓存；电脑负责较复杂且变化较快的应用逻辑；在线识别在独立工作进程内执行。这样的边界使传感器和通信实现能够被多个应用复用。

\fig{software}{主机软件的线程、进程与数据流}{GUI不直接打开第二个串口；识曲工作进程处理完整音频，游戏订阅已有数据消息。}

界面主线程处理用户操作、主题、页面与绘图；采集线程负责设备连接、轮询、命令排队和数据解码，再以消息形式发布结果。计算量较大或可能等待网络的识曲工作从界面分离，避免一次超时冻结整个窗口。SQLite负责歌曲、录音、历史、收藏及任务状态，原始文件通过任务编号与数据库关联。

\begin{table}[H]\centering\caption{软件模块与最小接口职责}\small
\begin{tabularx}{\linewidth}{p{28mm}Y Y}\toprule
模块 & 输入与输出 & 设计作用\\\midrule
设备与协议服务 & USB候选、命令队列$\to$状态/数据消息 & 持有唯一UART，检查帧、序号、构建及错误计数。\\
监测与事件控制 & 实时数据、用户配置$\to$曲线与完整事件 & 区分预览和快照，完成导出校验及波形回放。\\
音乐任务控制 & 三敲事件、录音状态$\to$完整WAV与识曲任务 & 组织录音、导出、匹配、收藏和中止。\\
识曲工作进程 & WAV、任务号与模式$\to$歌曲或失败原因 & 在线优先、本地回退；记录来源和耗时。\\
数据存储 & 任务、歌曲、结果$\to$SQLite与文件 & 历史可重复记录，收藏去重，文件与哈希关联。\\
游戏会话与面板 & 新鲜XYZ/按键$\to$校准、控制和存档 & 复用服务，管理互斥、暂停和配置恢复。\\\bottomrule
\end{tabularx}\end{table}

\section{启动、连接与状态恢复}
启动后先解析内置资源和持久化目录，加载主题与已有数据，再发现实际USB器件与串口。软件按设备序列号关联配置和UART通道，不把COM5等本机名称写成唯一端口。连接成功后读取身份、构建、传感器配置与缓冲状态；版本不匹配时提示初始化，避免用错误固件解释数据。

初始化将配套位流下载到FPGA SRAM。下载前确认板型与目标芯片，校验内置资源；常规FTDI和已有WinUSB两种后端用于适应驱动绑定差异。运行时不需要Vivado，但从源码重建硬件仍需要工具链。程序连接、重新初始化与板端复位形成不同运行阶段，记录各自版本和时间范围。

重连必须先读状态再决定动作。板上若存在冻结记录，优先保留并按原流程读取；掉电造成的SRAM数据丢失不会伪装成可恢复。对于游戏退出恢复，只有会话身份与配置条件仍成立时才恢复原参数；板端已复位时要求重新准备和校准，避免旧状态覆盖新会话。

\section{实时监测与事件导出流程}
监测线程周期获取板端快照和声音预览。包络、XYZ都使用序号去重，识别丢失区间；主机只将新数据追加到显示缓冲。暂停按钮冻结画面，后台采集仍继续。512点FFT快照按独立编号读取，使用原速样本计算频谱，不能从1 ms极值包络反推原始声音。

\fig{export}{完整事件读取与确认释放的顺序}{逐帧CRC、整事件CRC和文件哈希分别承担通信、完整性及归档追踪职责。失败读取可重试，确认释放发生在验证之后。}

事件导出先取得事件编号、长度、触发时刻、参数版本等元数据，再按偏移分块读取。每帧校验成功后检查返回编号与偏移；完整负载再核对板端总CRC、样本点数和时间连续性。通过后生成PCM、WAV、XYZ表和元数据，并发送确认释放。帧校验成功只能证明某个包未被检测到损坏，不能代替整条事件长度和内容一致性。

这种“采集冻结—编号读取—完整校验—确认释放”的事务式设计，把硬件临时状态转换为可重复使用的数据对象。用户可以回放录音、查看原始波形和导出文件；自动检测与手动抓取共用文件格式，有利于对照不同触发策略。

\section{听歌识曲的任务流程}
最终版通过SW15或界面进入就绪状态，连续三次符合间隔要求的敲击才启动。每两个相邻敲击的窗口约为0.15～1秒；一次或两次敲击不足以开始。就绪、录音、传输和匹配是不同阶段，忙碌时重复敲击不会创建并发录音。抢歌名玩法同样先以三敲启动，再记录人的拍桌锁定时刻。

\fig{music}{最终听歌识曲任务流程与故障回退}{默认自动模式优先在线；仅离线模式直接进入本地路径。中止针对当前任务，迟到结果通过任务身份过滤。}

默认20秒音乐录音完成后，主机先保存并核对完整WAV，再将中间12秒片段交给ShazamIO。完整文件和实际识别片段各有独立SHA-256。在线匹配成功后直接入库；未匹配、超时或服务不可用时尝试本地指纹。仅离线模式不发起在线请求，匹配依据不足则显示“未匹配”。

在线请求采用单次尝试和节流策略，请求间隔至少30秒；发生限流后冷却至少60秒，并遵从更长的服务端等待要求。冷却状态持久保存，防止重启绕开节流。识曲失败时原始录音仍可查看，用户明确重试时建立新任务，避免后台无界重试。

\section{历史、收藏与任务中止}
歌曲实体、每次识别历史和收藏是三个不同概念。同一首歌曲重复出现时，每次保留单独时间、来源与录音关联；收藏按歌曲去重。查询可按歌名、歌手、时间、来源和收藏状态筛选，次数统计在本机数据库完成。数据目录与程序内置资源分离，升级程序不要求覆盖用户历史。

任务中止需要同时处理板端和主机端。用户在录音、传输或匹配阶段发出中止后，服务向当前录音编号发送停止请求，并结束相关识曲工作；确认两端结束后才能重新就绪。若串口中止失败，界面保留“正在中止”或重试状态，不能立即把旧任务当作已停止。已保存的完整WAV、既有历史和收藏继续保留；已经发送到远端的请求无法撤回，但迟到结果不能覆盖新任务。

\section{体感游戏与平台协同}
《重力迷航》将加速度数据转为连续控制。输入模块对实际XYZ采样去重、平滑并检查新鲜度；以自然握持姿态建立相对零点，再确认方向映射。Z轴参与姿态有效性判断，不作为空间高度。250 ms没有新鲜XYZ时暂停，重复旧包不能延长有效期。

\fig{game}{体感游戏输入、逻辑与会话生命周期}{120 Hz为物理逻辑固定步长；实际绘制由主机调度。游戏订阅同一采集服务，不创建第二个UART线程。}

物理引擎采用120 Hz固定步长，使运动、碰撞与机关逻辑不随界面帧率直接变化；绘制复用Canvas图元，静态场景缓存后再进入运行。八关地图逐步引入路线、能量门、脉冲和移动障碍等规则，成绩、星级及会话记录独立保存。实板成绩与键盘练习分开，首次接入主程序时复制已有独立版进度，之后分别保存。

游戏进入前要求关闭SW15并退出识曲，暂停自动事件触发；游戏中暂缓完整事件导出与FFT快照，持续使用共享XYZ和按键。退出页面后恢复进入前的触发模式、开关和参数。页面离开、窗口失焦与最小化均暂停，返回后手动继续。该机制把多应用资源冲突处理纳入架构，而不是交由用户猜测哪个模块占用了设备。
